\chapter{Metodologija raziskave}\label{cha:metodologija}
V tem poglavju glede na tip naloge (raziskovalni ali razvojni) sistematično predstavite, razložite in utemeljite uporabljene \textbf{metode oz.~postopke} meritev, izračunov, modelirnih postopkov itn. ter predstavite in utemeljite izbor uporabljenih \textbf{materialov in vzorcev}. V tem poglavju posebej razdelajte tudi merilno negotovost, kjer je to smiselno. Metodologija mora biti opisana dovolj natančno, da je delo ponovljivo.

Vseh predlaganih podpoglavij ni obvezno vključiti v vsak tip naloge. Smiselno dodajte tudi druga podpoglavja.

\section{Preračuni}\label{sec:preracuni}
V tem podpoglavju so predstavljeni teoretični preračuni, analitični modeli ali numerični postopki, ki so podlaga za eksperimentalni ali razvojni del naloge.

Jasno navedite osnovne predpostavke modela, uporabljene enačbe ali izpeljave (z navedbo virov, če enačbe niso samostojno izpeljane), območje veljavnosti uporabljenih preračunov.

Po potrebi se v tem poglavju predstavi tudi numerični model (npr.~metoda končnih elementov), pri čemer se opišejo diskretizacija, robni pogoji in validacija modela.

\section{Eksperimentalni del}\label{sec:eksperimentalni_del}

\subsection{Vzorci in materiali}\label{sec:vzorci_materiali}
V tem podpoglavju je predstavljen izbor vzorcev in materialov. Namen je utemeljiti, zakaj so bili izbrani prav ti materiali in geometrije ter kako so povezani z namenom raziskave.

\subsection{Metodologija preizkusov}\label{sec:metodologija_preizkusov}
To podpoglavje opisuje eksperimentalni postopek in uporabljeno merilno opremo. Opis mora omogočati ponovljivost preizkusov.
